% Appendix for Section "hist" (writer: hist). Proofs adapted from research/tracks/history/notes-final.md
% (final numbering) and checked line by line. Steps the notes leave implicit are written out and marked by a
% "% step added" comment (not in the printed text). For the page cap, the secondary items lem:hist:memory,
% prop:hist:nopathwise, rem:hist:thm1open and lem:hist:chronolaw are stated here; prop:hist:normchain,
% prop:hist:strict, conj:hist:algmi, prop:hist:lhg, thm:hist:gf and thm:hist:expected were moved here from hist.tex
% in the arXiv edit pass (log-p1-hist-ex.md).
\section{Proofs for Section~\ref{sec:hist} (history-aware prediction)}\label{app:hist}

Notation is that of \cref{sec:hist}. We use: (F1) \cref{thm:setting:universal} for any countable decidable input set, $\MF(D)\ge2^{-K(T)}\nu_T(D)$ for $D\ne\emptyset$, and $\MF(D)\ge2^{-K(g)-c}$ if a partial computable $g$ extends $D$ (\cref{prop:setting:includes}(b)); (F2) the conditional coding theorem of \cref{sec:setting:data}: if $\nu(z\mid w)$ is lsc uniformly in $w$ and $\sum_z\nu(z\mid w)\le1$, then $K(z\mid w)\le-\log\nu(z\mid w)+K(\nu)+c$; (F3) Levin--Schnorr: $\omega$ is Martin-L\"of random iff $K(\omega_{1:\ell})\ge\ell-c_\omega$ for all $\ell$ \citep{li2008introduction}; (F4) $K(a,b)\le K(b)+K(a\mid b)+c$ and $K(s)\le|s|+K(|s|)+c$ \citep{li2008introduction}; (F5) $M\ge2^{-K(\nu)-c}\nu$ for every lsc semimeasure $\nu$ on $\{0,1\}^*$ \citep{zvonkin1970complexity,li2008introduction}; (F6) Sion's minimax theorem \citep{sion1958general}, for compact convex sets in $\mathbb R^m$ and a continuous payoff convex in one variable and affine in the other.

\subsection{The counterexamples}\label{app:hist:counter}

\begin{proof}[Proof of \cref{lem:hist:chrono}]
The events for different next labels are disjoint subsets of the event for $y_{1:n}$, and at a repeated input only the stored label survives; lower semicomputability is as for $\MF$. The second claim follows from \cref{cor:hist:H4}(a) with $K(\rho_F)=O(1)$. For the third, at a new input the chronological inequality is the weak inequality, and at a repeated input only one term is nonzero.
\end{proof}

\begin{lemma}[Chain semimeasure]\status{proved}\src{history Lemma 2.0}\label{lem:hist:chain0}
For a string $x$ and a sequence $z=(z_1,\dots,z_k)$ of nonempty strings ($k\ge0$) let $\nu(z\mid x):=\Unif\{\omega:f_\omega(x)=z_1,\,f_\omega(z_1)=z_2,\dots,f_\omega(z_k)=\epsilon\}$. Then $\sum_z\nu(z\mid x)\le1$, and $\nu$ is lsc uniformly in $x$.
\end{lemma}

\begin{proof}
Fix $\omega$, $x$, and $z\ne z'$ that both satisfy the condition; put $z_0=z'_0:=x$. If they first differ at $j\ge1$, then $z_j=f_\omega(z_{j-1})=z'_j$. If $z$ is a proper prefix of $z'$, of length $k$, then $\epsilon=f_\omega(z_k)=z'_{k+1}\ne\epsilon$. Both are contradictions, so the events are disjoint; they are effectively open, uniformly.
\end{proof}

\begin{proof}[Proof of \cref{thm:hist:chain}]
(a) Every $\omega$ with $f_\omega\supseteq D_n$ satisfies the condition of \cref{lem:hist:chain0} with $x=x_n$, $z=(x_{n-1},\dots,x_1)$, so $\MF(D_n)\le2^{-K(z\mid x_n)+c}$ by (F2). The program with $x_1,\dots,x_{n-1}$ built in that outputs $\epsilon$ at $x_1$, $x_{i-1}$ at $x_i$ ($2\le i\le n-1$) and $x_{n-1}$ elsewhere extends $D_n$; apply (F1).
(b) $x_{1:n}$ determines $\omega_{1:\bar\ell_n}$, so $K(x_{1:n})\ge\bar\ell_n-c_\omega$ by (F3); by (F4), $K(x_{1:n})\le\ell_n+K(\ell_n)+K(x_{1:n-1}\mid x_n)+c$. For the upper bound, $x_{1:n-1}$ is determined by $n$ and $\omega_{1:\bar\ell_{n-1}}$, and $\bar\ell_{n-1}$ is computable from $n$. For $\ell_i=i$ use $K(n)\le2\log n+c$.
(c) At step $i\ge2$ the history's inputs are $x_1,\dots,x_{i-1}$ and its labels $\epsilon,x_1,\dots,x_{i-2}$.
\end{proof}

\begin{proof}[Proof of \cref{thm:hist:skipped}]
(a) In block $\ell$ the strings other than $s_\ell$ come first, so $s_\ell$ is unseen when they are queried, and all its length-$\ell$ siblings are seen when it is. (b) $\dom D^{(\ell)}$ and $\ell$ are computable from each other, and given the domain the labelling determines $s_{1:\ell}$; apply \cref{thm:setting:coding}(a). A function machine that reads, for each length $j$, its own block of $j$ tape bits as a guess $s'_j$ and answers $1$ at $x$ iff $x=s'_{|x|}$ extends $D^{(\ell)}$ with probability $2^{-\ell(\ell+1)/2}$; apply (F1). (c) $K(s_{1:\ell}\mid\ell)\ge K(s_{1:\ell})-K(\ell)-c\ge\ell(\ell+1)/2-K(\ell)-c_\omega$ by (F3); and fewer than $2^{m-j}$ tuples have $K(\cdot\mid\ell)<m-j$, $m=\ell(\ell+1)/2$.
\end{proof}

\begin{proof}[Proof of \cref{prop:hist:orderfree}]
(a) For distinct nonempty $z_1,\dots,z_j\ne x$ let
\[O_j(z_{1:j}):=\{(x,z_1),(z_1,z_2),\dots,(z_{j-1},z_j)\},\qquad O_0:=\emptyset.\]
As $z_j$ (or $x$, if $j=0$) is not in $\dom O_j$, the weak inequality gives $\sum_y\nu(O_j\cup\{(z_j,y)\})\le\nu(O_j)$; the term $y=\epsilon$ is $\nu(D(x,z_{1:j}))$, each nonempty $y\notin\{x,z_1,\dots,z_j\}$ gives $\nu(O_{j+1}(z_{1:j},y))$, and the other terms are dropped. Induction on $J$ gives $\sum_{|z|\le J}\nu(D(x,z))\le\nu(\emptyset)\le1$.
(b) $D_n=D(x_n,(x_{n-1},\dots,x_1))$, and $z\mapsto\nu(D(x,z))$ is a semimeasure by (a), lsc uniformly in $x$; apply (F2) and continue as for \cref{thm:hist:chain}. For the skipped strings, the labellings of a finite set have total $\nu$-mass at most $1$, so $-\log\nu(D)\ge K(D\mid\dom D)-K(\nu)-c$.
(c) Gibbs' inequality given $x_n$ against the sub-probability $z\mapsto\nu(D(x_n,z))$; given $x_n$, $(x_{n-1},\dots,x_1)$ is uniform on $\bar\ell_{n-1}$ bits. For the skipped strings, against the labellings of the domain.
\end{proof}

\begin{lemma}[Normaliser lower bound]\status{proved}\src{history Lemma 2.4a}\label{lem:hist:broom}
For every $D$, $x\notin\dom D$ and function machine $B$: $\Znorm(D,x)\ge2^{-K(B)-c}\,\Unif\{\omega:f^B_\omega\supseteq D,\ f^B_\omega(x)\ne\bot\}$, which is $2^{-K(B)-c}\nu_B(D)$ if $f^B_\omega$ is always total.
\end{lemma}

\begin{proof}
Apply (F1) to each $D+(x,y)$ and sum over $y$; the events $\{f^B_\omega\supseteq D,\ f^B_\omega(x)=y\}$ are disjoint.
\end{proof}

For the chain, the \emph{broom} $B_j$ reads blocks $v_1,\dots,v_j$ of lengths $\ell_1,\dots,\ell_j$, sets $f(v_1):=\epsilon$, $f(v_i):=v_{i-1}$, and sends every other input to $v_j$; it extends $D_{j+1}$ with probability $2^{-\bar\ell_j}$ at a cost of $K(j)+c$ bits (as in \cref{prop:ex:normalised}).

\begin{proposition}[Normalised SOI on both examples]\status{proved; finite proxy computed}\src{history Prop 2.4, Lemma 2.4a; check c9}\label{prop:hist:normchain}
(a) In the setting of \cref{thm:hist:chain}(b), $-\log\MFbar(y_t\mid D_{t-1},x_t)\ge\ell_{t-1}-K(\ell_t)-K(t-2)-c_\omega$ for $t\ge2$; for $\ell_i=i$, $\LSFIbar\ge n(n-1)/2-n(4\log n+c_\omega)$. With independent uniform blocks the expected normalised loss at step $t$ is at least $\ell_{t-1}-K(t-2)-c$, so $\Ex\LSFIbar\ge n(n-1)/2-n(2\log n+c)$ for $\ell_i=i$.
(b) In the setting of \cref{thm:hist:skipped}(c), the normalised loss at $s_\ell$ is at least $\ell-2K(\ell)-c_\omega$, so $\LSFIbar\ge\ell(\ell+1)/2-\ell(4\log\ell+c_\omega)$ after block $\ell\ge2$; with independent uniform $s_j$, at least $\ell-K(\ell)-c$ in expectation.
\end{proposition}

\begin{proof}[Proof of \cref{prop:hist:normchain}]
Let $\beta_t:=-\log\MFbar(y_t\mid D_{t-1},x_t)=-\log\MF(D_t)+\log\Znorm(D_{t-1},x_t)\ge0$.
(a) The broom $B_j$ above is well defined (the $v_i$ have distinct lengths), total, and $K(B_j)\le K(j)+c$. On $\{v_i=x_i\ \forall i\le j\}$, of probability $2^{-\bar\ell_j}$, it extends $D_{j+1}$: $f(x_{j+1})=v_j=x_j$ since $\ell_{j+1}\notin\{\ell_1,\dots,\ell_j\}$. At step $t\ge2$, \cref{lem:hist:broom} with $B_{t-2}$ gives $\Znorm(D_{t-1},x_t)\ge2^{-K(t-2)-c-\bar\ell_{t-2}}$, so $\beta_t\ge-\log\MF(D_t)-\bar\ell_{t-2}-K(t-2)-c$. Insert \cref{thm:hist:chain}(b) at horizon $t$, or, for uniform blocks, \cref{prop:hist:orderfree}(c) in expectation. Sum the positive parts; for $\ell_i=i$ use $K(t),K(t-2)\le2\log n+c$ for $t\le n$.
(b) $T_\ell$ reads blocks $g_j$ of length $j$ for $j<\ell$ and outputs $1$ at $x$ iff $|x|<\ell$ and $x=g_{|x|}$, and $0$ elsewhere; it is total and $K(T_\ell)\le K(\ell)+c$. Before the query $s_\ell$ the data are $D':=D^{(\ell-1)}\cup\{(x,0):|x|=\ell,\,x\ne s_\ell\}$, which $T_\ell$ extends iff $g_j=s_j$ for all $j<\ell$, with probability $2^{-\ell(\ell-1)/2}$. So $\Znorm(D',s_\ell)\ge2^{-K(\ell)-c-\ell(\ell-1)/2}$; insert \cref{thm:hist:skipped}(c), or \cref{prop:hist:orderfree}(c) in expectation, and use $2K(j)\le4\log\ell+c$ for $j\le\ell$.
\end{proof}

The method needs more information per step than the $O(\log t)$ bits of the broom, so it gives nothing for \cref{prop:hist:strict,thm:hist:gf}.

\subsection{A dominance theorem for computable reveal laws}\label{app:hist:thm1}

\begin{lemma}[Private memory]\status{proved}\src{history Lemma 3.0}\label{lem:hist:memory}
Let a reveal mechanism that does not see $f$ consist of a random element $\Gamma$ (private memory and coins) independent of $f$ and maps $x_i=\mathrm{sel}_i(\Gamma,h_{i-1})$, with reveal-sequence law $\Delta(x_{1:n}\Vert y_{1:n-1}):=\Prob[\mathrm{sel}_i(\Gamma,h_{i-1})=x_i\ \forall i\le n]$. (a) Given $h_{i-1}$, $f$ and $x_i$ are conditionally independent; so $I(f\to x_{1:n})=0$ for finite-valued inputs and labels. (b) Given $f$, the input sequence has probability $\Delta(x_{1:n}\Vert f(x_1),\dots,f(x_{n-1}))$, and the effective kernel $\Delta(x_{1:i}\Vert y_{1:i-1})/\Delta(x_{1:i-1}\Vert y_{1:i-2})$ depends neither on $f$ nor on its law. (c) If the mechanism is a computable sampler with private memory that halts with probability $1$ at each step whatever labels it is fed, and its inputs up to time $n$ lie in a finite listed set $A_n$, then $\Delta(x_{1:n}\Vert y_{1:n-1})$ is a computable real, uniformly in its arguments.
\end{lemma}

\begin{proof}
The joint law of $(f,\Gamma,h_n)$ factorises (on countable spaces; in general via densities with respect to $\mathrm{law}(f)\otimes\mathrm{law}(\Gamma)$) as $[\mathrm{law}(f)\prod_i\ind[y_i=f(x_i)]]\cdot[\mathrm{law}(\Gamma)\prod_i\ind[x_i=\mathrm{sel}_i(\Gamma,h_{i-1})]]$. (a) Given $(h_{i-1},x_i)$ the law of $(f,\Gamma)$ is proportional to a function of $f$ times a function of $\Gamma$, and the $f$-factor does not involve $x_i$. (b) Sum the $\Gamma$-factor over $\Gamma$. (c) $\Delta(x\Vert y)$ is lsc (enumerate halting computations fed the labels $y$). Feed the sampler the labels $y_j$ while its inputs agree with $x_{1:j}$, and a fixed label otherwise; its input sequences in $A_n^n$ then have total probability $1$, and $x_{1:n}$ has probability $\Delta(x\Vert y)$. So $\Delta(x\Vert y)=1-\sum_{x'\ne x}\Prob[x']$, a finite sum of lsc terms subtracted from $1$.
\end{proof}

\begin{proof}[Proof of \cref{prop:hist:sketchfails}]
$P$ pays $1$ bit for $y_1$ and none for $y_2$. The simulator at $a$ outputs $P$'s uniform $y_1$; at $b$ it simulates $a$ then $b$ on fresh bits and outputs $y'_2=y'_1$, a uniform bit independent of the run at $a$; so $(\beta,\beta)$ gets probability $\frac14$. The mixture of the two constant functions gives it $\frac12$.
\end{proof}

\begin{proof}[Proof of \cref{thm:hist:shared}]
The function machine $T$ uses $\omega$ for one run of the alternating process: $x'_1\sim\Delta(\cdot\mid\emptyset)$, $y'_1\sim P(\cdot\mid\emptyset,x'_1)$, $x'_2\sim\Delta(\cdot\mid h'_1)$, \dots, each call reading fresh consecutive bits (continuing from the stored state if there is private memory). On query $x$ it runs until the first $j$ with $x'_j=x$ and outputs $y'_j$, diverging if there is none. Its answer depends on a finite prefix of $\omega$, so $T$ is a function machine, with $K(T)\le K(\Delta,P)+c$. If the trajectory begins with $x_1y_1\cdots x_ny_n$, each $x\in\dom D_n$ is first visited with its label in $D_n$ (the sequence is consistent), so $f^T_\omega\supseteq D_n$. The tapes producing this beginning are the disjoint union of the cylinders $[w_1v_1\cdots w_nv_n]$, $w_i$ ranging over the bits read by halting runs of $\Delta$ on $h_{i-1}$ with output $x_i$ and $v_i$ over those of $P$ on $(h_{i-1},x_i)$ with output $y_i$; each set is prefix-free, so the union has measure $\prod_i\Delta(x_i\mid h_{i-1})P(y_i\mid h_{i-1},x_i)$. % step added: measure of the union is the product
For distinct inputs, different orders give different trajectory prefixes of equal length, hence disjoint events. Apply (F1).
\end{proof}

\begin{corollary}[Deterministic reveal policies]\status{proved}\src{history Cor 3.3}\label{cor:hist:deterministic}
If $x_i=\delta(h_{i-1})$ for a computable $\delta$, then for every $f$, every sampler $P$ and every $n$, $\LSFI(D_n)\le L_P+K(\delta)+K(P)+c$ on every path.
\end{corollary}

\begin{corollary}[Against Solomonoff induction on the joint sequence]\status{proved}\src{history Cor 3.4}\label{cor:hist:joint}
With $M$ and $\code$ as in \cref{sec:hist:si}: $\LSFI(D_n)\le-\log M(\code(h_n))+c$ on every consistent sequence; for distinct inputs, $\LSFI(D_n)\le-\log\sum_\sigma M(\code(D_n\text{ in order }\sigma))+c$.
\end{corollary}

\begin{proof}[Proof of \cref{cor:hist:deterministic,cor:hist:joint}]
For \cref{cor:hist:deterministic}, take the input sampler $\delta$ in \cref{thm:hist:shared}; $L_\Delta=0$ on the realised path. For \cref{cor:hist:joint}, run the universal monotone machine on $\omega$, parse its output (uniquely) into blocks $\bar x'_1\bar y'_1\bar x'_2\cdots$, and answer $x$ with the label after the first input block equal to $\bar x$. If the output extends $\code(h_n)$ then $f_\omega\supseteq D_n$, and that event has probability $M(\code(h_n))$. For distinct inputs the codes of different orders have equal length and differ, so the events are disjoint.
\end{proof}

\begin{proof}[Proof of \cref{thm:hist:bayes}]
The joint law of $(f,\Gamma,h_n)$ is $\nu(df)\Prob(d\Gamma)\prod_i\ind[x_i=\mathrm{sel}_i(\Gamma,h_{i-1})]\ind[f\supseteq D_n]$, a function of $f$ times a function of $\Gamma$ for fixed $h_{i-1}$. So given $h_{i-1}$ and $x_i$ the law of $f$ is $\nu(\cdot\mid f\supseteq D_{i-1})$, whose prediction for $y_i$ is the Bayes predictor; the product telescopes to $\nu(D_n)$, and (F1) gives $\MF\ge2^{-K(\nu)}\nu$. Gibbs' inequality, valid for sub-probabilities, gives $\Ex L_P\ge\Ex[-\log\nu(D_n)]$ step by step.
\end{proof}

\begin{proof}[Proof of \cref{thm:hist:directed}]
(a) Let $Q_f$ be a regular conditional law of $x_{1:n}$ given $f$. Then $\Prob(h_n\mid f)=Q_f(x_{1:n})\ind[f\supseteq D_n]$ and $\Prob(h_n)=\nu(D_n)\bar Q_D(x_{1:n})$ with $\bar Q_D:=\Ex_{f'\sim\nu}[Q_{f'}\mid f'\supseteq D_n]$; the posterior of $f$ has density $Q_f/\bar Q_D$ with respect to $\nu(\cdot\mid f\supseteq D_n)$. Hence $I(f;h_n)=\Ex\KL(\Prob(f\in\cdot\mid h_n)\Vert\nu(\cdot\mid f\supseteq D_n))+\Ex[-\log\nu(D_n)]$. By the chain rule $I(f;h_n)=\sum_i[I(f;x_i\mid h_{i-1})+I(f;y_i\mid h_{i-1},x_i)]=I(f\to x)+\Ex L_\Prob$, because $y_i$ is a function of $(f,x_i)$. All terms are finite ($I(f;h_n)\le H(h_n)$). (b) (F1), then (a), then $\Ex L_P\ge\Ex L_\Prob$ by Gibbs' inequality.
\end{proof}

\begin{proof}[Proof of \cref{thm:hist:randchain}]
The machine compares a query of length $i$ with tape block $i$ and returns block $i-1$ if they agree and $i\ge2$, and $\epsilon$ otherwise. Given $h_{i-1}$, $x_i$ is uniform on $\{0,1\}^i$; given $f$ it is determined; so $I(f;x_i\mid h_{i-1})=i$. The losses are \cref{prop:hist:orderfree}(c) and \cref{prop:hist:normchain}(a).
\end{proof}

\begin{proof}[Proof of \cref{thm:hist:minimax}]
Write $Y^A:=Y^{A_n}$, $\mathcal U$ for the uniform law on it, and $\mathcal G:=\{\frac12G_0+\frac12\mathcal U\}$ over laws $G_0$ on $Y^A$; $\mathcal G$ is compact and convex, and $G(D)\ge\frac12|Y|^{-n}$ when $|D|\le n$. Let $\mathcal F_{\rm fin}:=\{f:\Ex_{x\mid f}L_P<\infty\}$ (other $f$ are trivial), $\Phi_f(G):=\Ex_{x\mid f}[-\log G(D_n)]-\Ex_{x\mid f}L_P$, and $\Phi(G,\nu):=\sum_f\nu(f)\Phi_f(G)$ for laws $\nu$ on $\mathcal F_{\rm fin}$.
\emph{Step 1.} $\Phi$ is continuous, convex in $G$, affine in $\nu$; by (F6), $\min_G\max_\nu\Phi=\max_\nu\min_G\Phi$. For $G=\frac12\nu+\frac12\mathcal U$, $-\log G(D)\le1-\log\nu(D)$, so \cref{thm:hist:directed}(a) and Gibbs' inequality give $\Phi(G,\nu)\le1+I_\nu(f\to x)$. Hence $v^*:=\min_G\max_{f\in\mathcal F_{\rm fin}}\Phi_f(G)\le1+C_{\rm dir}$.
\emph{Step 2.} For $G,G'\in\mathcal G$ and $|D|\le n$, $|\log G(D)-\log G'(D)|\le|G(D)-G'(D)|/(\ln2\min(G(D),G'(D)))\le|Y|^n\|G_0-G'_0\|_1/\ln2$. So a rational $G'_0$ within $\ell^1$-distance $|Y|^{-n}\ln2$ of a minimiser has $\max_{f\in\mathcal F_{\rm fin}}\Phi_f(G')\le v^*+1\le2+C_{\rm dir}$.
\emph{Step 3, case (a).} Path probabilities given $f$ are computable rationals, $\mathcal F_{\rm fin}$ is decidable ($f\notin\mathcal F_{\rm fin}$ iff some path of positive probability has $P(y_i\mid h,x)=0$), and each $\Phi_f$ is computable on rational points. With the Lipschitz constant of Step 2, evaluating $\max_f\Phi_f$ to within $\frac14$ on a fine rational grid yields a grid point $G'$ with $\max_f\Phi_f(G')\le v^*+1$.
\emph{Step 3, case (b).} Let $\gamma$ be rational with $C_{\rm dir}\le\gamma\le C_{\rm dir}+1$ ($\gamma=0$ if $\Delta$ does not see $f$). Enumerate rational $G'_0$ and stop at the first $G'$ certified to satisfy $(\star)$: $\Ex_{x\mid f}[-\log G'(D_n)]<\Ex_{x\mid f}L_P+3+\gamma$ for every $f\in Y^A$. The left side is a computable real. The right side is a finite sum of products of a computable nonnegative real $\Prob[\text{path}\mid f]$ and an lsc value $L_P(\text{path})\in[0,\infty]$, hence lsc in $[0,\infty]$ (with $0\cdot\infty=0$). A true strict inequality ``computable $<$ lsc'' is certified in finite time, also when the right side is $+\infty$; dovetail over candidates. By Step 2 some rational $G'$ satisfies $(\star)$, so the search halts. If $\Delta$ sees $f$, $\gamma$ must be computed; this case is only sketched (\cref{rem:hist:minimaxsketch}).
\emph{Step 4.} The procedure is uniform in $(P,\Delta,n,A)$, so $G'$ has rational probabilities computable from these, and rejection sampling (undefined outside $A_n$) makes it a function machine with $K(G')\le K(P)+K(\Delta)+K(n)+K(A)+c$. (F1) gives $\Ex_{x\mid f}\LSFI\le\Ex_{x\mid f}[-\log G'(D_n)]+K(G')+c\le\Ex_{x\mid f}L_P+4+C_{\rm dir}+K(G')+c$; average over $f$ for the last claim. An almost surely halting reveal sampler with private memory has $K(\Delta)$ at most its complexity plus $K(A)+c$ (\cref{lem:hist:memory}(c) uses the list $A_n$).
\end{proof}

\begin{proof}[Proof of \cref{cor:hist:thm1}]
(a) \Cref{thm:hist:minimax}(b) with $\Delta:=\Delta_f$, which does not see $f$; then average. (b) \Cref{cor:hist:deterministic}. (c) Take expectations in \cref{thm:hist:shared}; given $f$ the law of $x_{1:n}$ is $\Delta_f(x\Vert y)$, and $\Ex_{x\mid f}[-\log\Delta_f(x\Vert y)]=H(x_{1:n}\mid f)$, as $\Delta_f$ halts with probability $1$.
\end{proof}

\begin{remark}[Open cases in Formulation C]\status{open}\src{history remarks after Cor 3.9; Q1, Q7; checks r6, r7, c3, r2}\label{rem:hist:thm1open}
(i) Whether $K(n)+K(A)$ can be dropped: the minimax law depends on the horizon; in one small instance a direct computation (code in the accompanying repository) certifies that one law for horizons $1$--$4$ needs at least $0.026$ bits, while each single horizon needs at most $0.005$, and it finds no growth up to horizon $6$. (ii) Defective samplers with $K(P)$ alone ($\Ex L_P$ is then only upper semicomputable, so $(\star)$ cannot be certified). (iii) Infinite alphabets at a fixed horizon. (iv) Reveal laws that see $f$, in case (b) (\cref{rem:hist:minimaxsketch}).
\end{remark}

\begin{proposition}[No pathwise version of Formulation A]\status{proved}\src{history Prop 3.12}\label{prop:hist:nopathwise}
No bound $\LSFI(D_n)\le L_P+K(P)+K(\Delta)+c$ holds on every path for every fixed $f$ and every reveal law $\Delta$ that does not see $f$.
\end{proposition}

\begin{proof}
Let $\omega$ be Martin-L\"of random with blocks $b_i$ of length $i$, $f(b_{i+1}):=b_i$ and $f:=\epsilon$ elsewhere; let $\Delta$ draw independent uniform strings of lengths $1,2,\dots$, and $P$ be the chain predictor. The path $x_i=b_i$ has probability $2^{-n(n+1)/2}>0$; on it $L_P=0$ and $\LSFI(D_n)\ge n(n-1)/2-2\log n-c_\omega$ (\cref{thm:hist:chain}(b)).
\end{proof}

\begin{proposition}[Shannon form of the mutual-information bound]\status{proved; computed}\src{history Prop 3.10; check c7}\label{prop:hist:shannon}
Let $(f,\Theta)$ be jointly random and, given $(f,\Theta)$, let $x_i\sim\Delta_{\Theta,\Gamma}(\cdot\mid h_{i-1})$, a kernel that does not see $f$, with private memory $\Gamma$ independent of $f$ given $\Theta$; finite alphabets up to time $n$. (a) If $I(f;\Theta)<\infty$, then $I(f\to x_{1:n})=I(f;\Theta)-I(f;\Theta\mid h_n)-\sum_iI(y_i;\Theta\mid h_{i-1},x_i)\le I(f;\Theta)$. (b) For $f\sim\nu$ computable, $\Ex\LSFI(D_n)\le\Ex L_P+I(f;\Theta)+K(\nu)+c$ for every $P$. (c) In Formulation C, $\Theta:=\Delta^{(n)}_f$, the restriction of $\Delta_f$ to $n$ steps, is a function of $f$, $I(f;\Theta)=H(\Theta)$, and in \cref{thm:hist:randchain} $I(f\to x)=H(\Theta)=n(n+1)/2$.
\end{proposition}

\begin{proof}
(a) Given $\Theta=k$, the law of $(f,\Gamma,h_{i-1})$ is $p(df\mid k)p(d\Gamma\mid k)\prod_{j<i}\Delta_{k,\Gamma}(x_j\mid h_{j-1})\ind[f(x_j)=y_j]$, so $f\perp\Gamma\mid(\Theta,h_{i-1})$, and the law of $x_i$ given $(f,\Theta,h_{i-1})$ does not involve $f$: $I(f;x_i\mid\Theta,h_{i-1})=0$. The chain rule gives $I(f;\Theta,h_n)=I(f;\Theta)+\sum_iH(y_i\mid\Theta,h_{i-1},x_i)$, and also $I(f;\Theta,h_n)=I(f;h_n)+I(f;\Theta\mid h_n)$ with $I(f;h_n)=I(f\to x)+\sum_iH(y_i\mid h_{i-1},x_i)$ (proof of \cref{thm:hist:directed}). Equate, using $H(y_i\mid h_{i-1},x_i)-H(y_i\mid\Theta,h_{i-1},x_i)=I(y_i;\Theta\mid h_{i-1},x_i)$. (b) \Cref{thm:hist:directed}(b) and (a). (c) $I(f;\Theta)=H(\Theta)$ when $\Theta$ is a function of $f$; in the randomised chain, $\Delta^{(n)}_f$ determines the chain, uniform on $n(n+1)/2$ bits. % restriction to n steps added (the notes write H(Delta_f))
\end{proof}

\begin{conjecture}[Algorithmic form in Formulation C]\status{conjecture (untested)}\src{history Conj 3.11; Q8}\label{conj:hist:algmi}
With $I(a:b):=K(b)-K(b\mid a)$ and the hypotheses of \cref{cor:hist:thm1}(a): $\Ex_{x\mid f}\LSFI\le\Ex_{x\mid f}L_P+I(\Delta_f:f|_{A_n})+K(P)+O(\log K(\Delta_f))+K(n)+K(A)+c$.
\end{conjecture}

We have neither a proof nor a counterexample.

\subsection{\SFIhist}\label{app:hist:sfihist}

\begin{lemma}\status{proved}\src{history Lemma 4.1}\label{lem:hist:chronolaw}
$M^{\rm hist}_T$ is an lsc chronological semimeasure, uniformly in an index of $T$.
\end{lemma}

\begin{proof}
The events $\{T^\omega(x_{1:i})=y_i\ \forall i\le n,\ T^\omega(x_{1:n}x)=y\}$ for different $y$ are disjoint subsets of the event for $(y_{1:n},x_{1:n})$ and are effectively open, uniformly (\cref{prop:setting:machinelaw}); the label at time $i$ depends only on $x_{1:i}$.
\end{proof}

\begin{proof}[Proof of \cref{thm:hist:representation}]
Nodes are $z=(x_{1:k},y_{1:k})$ with parent $p(z)$; the \emph{family} $F(z)$ of a non-root node consists of the nodes with the same parent and the same last input. Take dyadic, finitely supported $\rho_t\uparrow\rho$, computable and nondecreasing in $t$, and a computable enumeration $z_1,z_2,\dots$ of the non-root nodes with every node after its parent. Each node carries a finite union $S(z)$ of dyadic intervals (cylinders), initially empty; $S(\mathrm{root}):=[0,1)$. At stage $t$ process $z_1,\dots,z_t$ in order: with $\mathrm{free}(z):=S(p(z))\setminus\bigcup_{z'\in F(z)}S(z')$ (current sets, $z$ included), if $a:=\min(\rho_t(z)-\Unif(S(z)),\Unif(\mathrm{free}(z)))>0$, add to $S(z)$ a dyadic part of $\mathrm{free}(z)$ of measure $a$. By construction: (I1) $S(z)\subseteq S(p(z))$; (I2) family members are disjoint; (I3) $\Unif(S_t(z))\le\rho_t(z)$ after stage $t$; (I4) $S_t(z)$ grows with $t$.
\emph{Limit.} We show $\Unif(S(z))=\rho(z)$, $S(z):=\bigcup_tS_t(z)$, by induction along the enumeration. With $p:=p(z)$, the hypothesis (or $\Unif(S(\mathrm{root}))=1$) and the chronological inequality give $\Unif(S(p))\ge\sum_{z'\in F(z)}\rho(z')$. When $z$ is processed at stage $t$, $S(p)=S_t(p)$ ($p$ is processed earlier), and each sibling has measure at most $\rho(z')$ by (I3). So afterwards $\Unif(S_t(z))\ge\min(\rho_t(z),\Unif(S_t(p))-\sum_{z'\in F(z)\setminus\{z\}}\rho(z'))$; only ``$\ge$'' holds, since a later sibling may still grow. Letting $t\to\infty$ gives $\liminf_t\Unif(S_t(z))\ge\rho(z)$, and (I3) gives equality.
\emph{Machine.} $T^\omega_\rho(x_{1:k})$ recomputes $y_i:=T^\omega_\rho(x_{1:i})$ for $i<k$, then searches for $t,y$ with $\omega\in S_t((x_{1:k},y_{1:k-1}y))$ and outputs the first $y$ found; by (I2) at most one $y$ is ever found, and membership depends on a finite prefix of $\omega$. By induction with (I1), (I2), $\{\omega:T^\omega_\rho(x_{1:i})=y_i\ \forall i\le n\}=S((x_{1:n},y_{1:n}))$, of measure $\rho(y_{1:n}\Vert x_{1:n})$. The construction is uniform in an index of $(\rho_t)$.
\end{proof}

\begin{proof}[Proof of \cref{cor:hist:H4}]
(a) \Cref{thm:hist:representation} and (F1) over $X_H$ at non-root nodes; at the root $\MH=1\ge\rho(\epsilon\Vert\epsilon)$; two universal elements dominate each other. (b) $P(y\Vert x)$ is an lsc chronological semimeasure of complexity at most $K(P)+c$. (c) \Cref{lem:hist:chrono}.
\end{proof}

\begin{proposition}[Other oracles]\status{proved}\src{history Prop 4.4}\label{prop:hist:variants}
Let the inputs be distinct. (a) $\Mpairs$ and $\MH$ agree up to constant factors. (b) \SFIcount{} is SOI on stamped inputs $(i,x_i)$ (as in \cref{thm:learn:batch}(b)), and $L_{\rm count}\le-\log M(\bar y_1\cdots\bar y_n)+c$. (c) SOI $\preceq$ \SFIcount{} $\preceq$ \SFIset{} $\preceq$ \SFIhist{}, where $\nu\preceq\nu'$ means $\nu\le2^c\nu'$.
\end{proposition}

\begin{proof}
(a) A pairs-machine $T$ gives the history-oracle machine
\[T'^\omega(x_{1:i}):=T^\omega(x_1,T'^\omega(x_1),\dots,x_{i-1},T'^\omega(x_{<i}),x_i);\]
on the event that $T'$ outputs $y_{1:n}$ the recursion reproduces the true history, and conversely by induction on $i$. (b) The hypothesis that outputs the $i$-th $\bar y$-block of the universal monotone machine's output. (c) Each oracle is computable from the next (the count is the size of the set plus one).
\end{proof}

\begin{proposition}[The hierarchy is strict]\status{proved; lemma computed}\src{history Prop 4.5; checks c6, r4}\label{prop:hist:strict}
Semimeasure convention. (i) Every set $S$ of $n\ge2$ strings has an order such that, with labels $y_i=\ind[i\text{ even}]$, $\LSFI\ge n-\log(n+1)-c$ while $L_{\rm count}\le c$. (ii) On the chain of \cref{thm:hist:chain}(b), $L_{\rm count}\ge\bar\ell_{n-1}-K(\ell_n)-c_\omega$ while $L_{\rm set}\le c$. (iii) Every such $S$ has an order such that, with labels $y_1=\epsilon$, $y_i=x_{i-1}$, $L_{\rm set}\ge\log n!-\log n-c$ while $\LH\le c$. Single infinite sequences with linearly growing gaps also exist (\cref{prop:hist:strictinf}).
\end{proposition}

\begin{proof}[Proof of \cref{prop:hist:strict}]
(i) $A:=\{x_i:i\text{ even}\}$ and $D$ are computable from each other given $S$, so $\LSFI\ge K(A\mid S)-c$ (\cref{thm:setting:coding}(a)); choose $A$ of size $\lfloor n/2\rfloor$ with $K(A\mid S)\ge\log\binom n{\lfloor n/2\rfloor}-1\ge n-\log(n+1)-1$. The count hypothesis $\ind[i\text{ even}]$ gives $L_{\rm count}\le c$.
(ii) With $U$ the universal count-machine, $\nu(z\mid x,n):=\Unif\{\omega:U^\omega(n,x)=z_1,\,U^\omega(n-1,z_1)=z_2,\dots,U^\omega(1,z_{n-1})=\epsilon\}$ sums to at most $1$ over $z$ (iteration determines $z$), so $L_{\rm count}\ge K(x_{1:n-1}\mid x_n,n)-c$; $n$ is computable from $|x_n|$, and the proof of \cref{thm:hist:chain}(b) finishes. ``Output the longest element of the set, or $\epsilon$ if it is empty'' gives $L_{\rm set}\le c$.
(iii) For a tape $\omega$ of the universal set-machine $U$, an order is consistent iff $U^\omega(\emptyset,x_1)=\epsilon$ and $U^\omega(\{x_{<i}\},x_i)=x_{i-1}$. Then $x_n$ determines $x_{n-1}=U^\omega(S\setminus\{x_n\},x_n)$, and so on backwards, so at most $n$ orders are consistent with $\omega$ and $\sum_\sigma\Mset(\sigma)\le n$. (F2) for this semimeasure divided by $n$, lsc uniformly in $S$, gives $L_{\rm set}\ge K(\sigma\mid S)-\log n-c$; choose $K(\sigma\mid S)\ge\log n!-1$. ``$y_1=\epsilon$, $y_i=x_{i-1}$'' gives $\LH\le c$.
\end{proof}

\begin{proposition}[Single infinite sequences]\status{proved}\src{history Prop 4.5$'$}\label{prop:hist:strictinf}
Let $\omega$ be Martin-L\"of random and $u_j,v_j$ the strings with indices $2j-1,2j$ in length-lexicographic order; block $j$ reveals $u_j,v_j$ in this order if $\omega_j=0$, reversed otherwise. (i) With labels $\ind[i\text{ even}]$, $\LSFI(D_{2k})\ge k-K(k)-c_\omega$ for all $k$, while $L_{\rm count}\le c$. (ii) The chain of \cref{thm:hist:chain}(b) is a single sequence for \cref{prop:hist:strict}(ii). (iii) With labels $y_1=\epsilon$, $y_i=x_{i-1}$, $L_{\rm set}\ge k-K(k)-c_\omega$ at time $2k$, while $\LH\le c$. So each gap grows linearly along one sequence (semimeasure convention).
\end{proposition}

\begin{proof}
(i) $\dom D_{2k}$ is computable from $k$, and the labelling determines $\omega_{1:k}$; use \cref{thm:setting:coding}(a), (F4) and (F3). (iii) By the backward reconstruction, a tape fixes the order given its last input ($u_k$ or $v_k$), so at most two of the $2^k$ block orders $\sigma_b$ fit one tape; apply (F2) to $b\mapsto\Mset(\sigma_b)/2$ and take $b=\omega_{1:k}$.
\end{proof}

$\Mpq$ is Solomonoff's per-query solution, his eq.~(10) for the operator-induction problem (\cref{rem:versions:operator}).

\begin{proposition}[Solomonoff's per-query solution ignores the history]\status{proved}\src{history Prop 4.9}\label{prop:hist:operator}
$\Mpq$ is an lsc chronological semimeasure that ignores the history, so $\LH\le-\log\Mpq+c$ on every sequence. On distinct inputs $\Mpq=\Mprod(D_n)$, a weak semimeasure, so it pays the lower bounds of \cref{thm:hist:chain,thm:hist:skipped} (semimeasure convention).
\end{proposition}

\begin{proof}
\Cref{prop:learn:perquery}(a), \cref{cor:hist:H4}(a), and $\Mprod\le2^c\MF$ (\cref{thm:versions:chain}); $\Mprod$ is weak, so \cref{prop:hist:orderfree}(b) applies.
\end{proof}

\begin{proposition}[Deterministic programs for \SFIhist]\status{proved; toy computed}\src{history Prop 4.10; checks s2, c10}\label{prop:hist:programs}
With $(\varphi_e)$ the partial computable functions on $X_H$, let $\Whist(y\Vert x):=\sum_{p\in\dom V}2^{-|p|}\ind[\varphi_{V(p)}(x_{1:i})=y_i\ \forall i\le n]$ and $\Khist(y\Vert x)$ the least such $|p|$; let $\Ktail(n):=-\log\sum_{N\ge n}2^{-K(N)}$ and $\{0,1\}\subseteq Y$. For every input list and $y\in\{0,1\}^n$, $\Whist(y\Vert x)\ge2^{-n-\Ktail(n)-c}$, where (a) $\Ktail(n)\le\Kpad(n)\le K(n)$, (b) $\Ktail(n)\le g(n)+K(g)+c$ for every computable nondecreasing unbounded $g$, and (c) $\Ktail(n)\to\infty$. (d) But $n-c\le\max_y\Khist(y\Vert x_{1:n})\le n+K(n)+c$, and along every input sequence $\max_y\Khist(y\Vert x_{1:n})-n>\log n+\log\log n$ for infinitely many $n$.
\end{proposition}

\begin{proof}
$\Whist$ is an lsc chronological semimeasure (programs consistent with different next labels differ). Let $\varphi_{N,w}(x_{1:i}):=w_i$ for $i\le N$, $w\in\{0,1\}^N$. A prefix machine reads a $V$-program for $N$ and then $N$ bits $w$ and outputs an index of $\varphi_{N,w}$; by universality, with a fixed string $r$ and $q_N$ a shortest program for $N$, the strings $rq_Nw$ are distinct elements of $\dom V$ ($q_N,q_{N'}$ are incomparable for $N\ne N'$), of length $|r|+K(N)+N$. For $N\ge n$, $\varphi_{N,w}$ is consistent iff $w_{1:n}=y$, for $2^{N-n}$ strings $w$, so $\Whist(y\Vert x)\ge\sum_{N\ge n}2^{N-n-|r|-K(N)-N}=2^{-n-|r|-\Ktail(n)}$.
(a) One term of the sum. (b) $q(N):=2^{-g(N)}-2^{-g(N+1)}$ is a computable discrete measure with $\sum_{N\ge n}q(N)=2^{-g(n)}$, and (F2) gives $2^{-K(N)}\ge2^{-K(g)-c}q(N)$. (c) The tails of $\sum_N2^{-K(N)}\le1$ tend to $0$.
(d) The upper bound is the program $rq_ny$; the lower bound holds since a consistent program and $x_{1:n}$ determine $y$. Fix the input sequence, let $e(n):=\max_y\Khist(y\Vert x_{1:n})-n$, $e^*(n):=\max_{m\le n}e(m)$, and let $C(L)$ and $k_L$ count the $p\in\dom V$ with $|p|\le L$ and $|p|=L$. Distinct labellings need distinct programs, so $C(n+e^*(n))\ge2^n$. By summation by parts and Kraft, $\sum_{L\le m}C(L)2^{-L-1}\le\sum_Lk_L2^{-L}\le1$; as $n+e^*(n)$ is strictly increasing, $\sum_n2^{-e^*(n)}\le2$. If $e(n)\le\log n+\log\log n$ for all $n\ge n_0$, then $\sum_n2^{-e^*(n)}\ge2^{-O(1)}\sum_n1/(n\log n)=\infty$.
\end{proof}

\subsection{Solomonoff induction on the joint sequence}\label{app:hist:si}

Block sequences are $b=(x_1,y_1,x_2,\dots)$, possibly ending after an input; $b^-$ is $b$ without its last block.

\begin{lemma}[Block codes]\status{proved}\src{history Lemma B; check r4}\label{lem:hist:block}
Let $\nu$ be an lsc semimeasure on block sequences ($\nu(\emptyset)\le1$, $\sum_\beta\nu(b\beta)\le\nu(b)$), and $s\mapsto\bar s$ a prefix-free code with nonempty codewords. Let $B(z)$ be the set of $b$ with $\code(b)\sqsupseteq z$ and ($b=\emptyset$ or $\code(b^-)\not\sqsupseteq z$). Then $\nu_{\rm bin}(z):=\sum_{b\in B(z)}\nu(b)$ is an lsc semimeasure on $\{0,1\}^*$ with $\nu_{\rm bin}(\code(b))=\nu(b)$, so $M(\code(b))\ge2^{-K(\nu)-c}\nu(b)$.
\end{lemma}

\begin{proof}
$B(\code(b))=\{b\}$: codes are prefix-free and nonempty, so $\code(b')\sqsupseteq\code(b)$ forces $b'$ to extend $b$ blockwise, and if $b'\ne b$ then $\code(b'^-)\sqsupseteq\code(b)$. $B(\epsilon)=\{\emptyset\}$. Let $b\in B(za)$. Either $\code(b^-)\not\sqsupseteq z$, and $b\in B(z)$ with $\code(b)\ne z$; or $\code(b^-)\sqsupseteq z$, and since $\code(b^-)$ and $za$ are prefixes of $\code(b)$ with $\code(b^-)\not\sqsupseteq za$, $\code(b^-)=z$, so $b$ extends by one block the element $b^-$ of $B(z)$ with code $z$. Conversely each $b\in B(z)$ with $\code(b)\ne z$ lies in exactly one $B(za)$. Hence $\nu_{\rm bin}(z0)+\nu_{\rm bin}(z1)\le\sum_{b\in B(z),\code(b)\ne z}\nu(b)+\sum_{b'\in B(z),\code(b')=z}\sum_\beta\nu(b'\beta)\le\nu_{\rm bin}(z)$. Apply (F5).
\end{proof}

\begin{proof}[Proof of \cref{thm:hist:si}]
(a) $M(\code(h_n))$ is an lsc chronological semimeasure: with $z:=\code(h_n)$, $\sum_yM(z\bar x\bar y)\le M(z\bar x)\le M(z)$, as $\{\bar y\}$ is prefix-free. Apply \cref{cor:hist:H4}(a). Likewise $M(\code(h_{n-1})\bar x_n)\ge M(\code(h_n))$ is one for inputs given labels. (b) $\nu(h_n):=\MHin(x_{1:n}\Vert y_{1:n-1})\MH(y_{1:n}\Vert x_{1:n})$ and $\nu(h_{n-1}x_n):=\MHin(x_{1:n}\Vert y_{1:n-1})\MH(y_{1:n-1}\Vert x_{1:n-1})$ define an lsc semimeasure on block sequences (sum over $y_n$ with $\MH$'s chronological inequality, over $x_n$ with $\MHin$'s); apply \cref{lem:hist:block}. (c) Combine with $\LMjoint=\LMx+\LMy$. (d) \Cref{prop:hist:inputloss}.
\end{proof}

\begin{proof}[Proof of \cref{prop:hist:inputloss}]
$\nu(h_n):=\ind[x_i=\delta(h_{i-1})\ \forall i\le n]\MH(y_{1:n}\Vert x_{1:n})$ and $\nu(h_{n-1}x_n):=\ind[x_i=\delta(h_{i-1})\ \forall i\le n]\MH(y_{1:n-1}\Vert x_{1:n-1})$ define an lsc semimeasure on block sequences with $K(\nu)\le K(\delta)+c$: the indicator is c.e., summing over the next input leaves only $\delta(h_{n-1})$, and summing over the next label uses $\MH$'s chronological inequality. \Cref{lem:hist:block} gives $\LMjoint\le\LH+K(\delta)+c$; with $\LMjoint=\LMx+\LMy$ this is the left inequality. The right one is \cref{thm:hist:si}(c).
\end{proof}

On $\delta$-consistent sequences also $L^x_{\rm in}\le K(\delta)+c$, so $\LH+L^x_{\rm in}-K(\delta)-c\le\LMjoint\le\LH+L^x_{\rm in}+c$.

Lattimore, Hutter and Gavane (LHG below) \citep{lattimore2011universal} show that unnormalised Solomonoff induction can fail to predict simple recursive sub-patterns of a sequence. We do not show that their construction gives the following assumptions, which are not proved here.

\begin{assumption}[LHG-even]\label{ass:hist:lhgeven}
With $X=Y=\{0,1\}$ and the identity code, for some $\omega$ the loss of unnormalised $M$ on positions $2,4,\dots$ of $\omega_1\omega_1\omega_2\omega_2\cdots$ is unbounded.
\end{assumption}

\begin{assumption}[LHG-odd]\label{ass:hist:lhgodd}
The same for positions $3,5,\dots$ of $0\omega_1\omega_1\omega_2\omega_2\cdots$.
\end{assumption}

\begin{proposition}[The conditional label loss]\status{conditional (on LHG transfer); LHG known, unverified beyond the abstract}\src{history \S4.4, Q3}\label{prop:hist:lhg}
With $X=Y=\{0,1\}$ and the identity code: under Assumption LHG-even, $\LMy\le\LH+O(1)$ fails on some sequence (labels copy inputs); under Assumption LHG-odd, $\LH\le\LMy+O(1)$ fails on some sequence (inputs copy the previous label, \cref{prop:hist:inputloss}). For the code $\bar s=1^{|s|}0s$ the assumptions must also hold for framed blocks.
\end{proposition}

\begin{proof}[Proof of \cref{prop:hist:lhg}]
(Even.) $x_i:=\omega_i$, $y_i:=x_i$: $\LH\le c$ by the copy predictor (\cref{cor:hist:H4}(b)), while $\LMy$ is $M$'s loss on the copied bits. (Odd.) $\delta(\emptyset):=0$, $\delta(h):=$ the last label of $h$, $y_i:=\omega_i$: the joint sequence is $0\omega_1\omega_1\omega_2\omega_2\cdots$, the inputs $x_2,x_3,\dots$ are the copied bits, so $\LMx$ is unbounded, and \cref{prop:hist:inputloss} with $K(\delta)=O(1)$ gives $\LH-\LMy\to\infty$. Non-convergence of the copied bits' conditional probabilities to $1$ already makes the cumulative loss unbounded.
\end{proof}

\subsection{What SOI can and cannot exploit}\label{app:hist:exploit}

\begin{proof}[Proof of \cref{thm:hist:sandwich,cor:hist:initial}]
The left inequality is \cref{lem:hist:chrono}. For the right one fix $\sigma'$ and $p$ attaining $\Kpred(x_{\sigma'})$. The function machine $T$, which does not depend on $p$, on query $x$ enumerates the prefix-free set $\dom V(\cdot,x)$ and takes the first element $q$ found that is a prefix of the tape; checking a candidate reads the tape only up to the first mismatch, which comes within $|q|$ bits, so $T$ reads exactly $q$. % step added: T reads exactly q
It computes $V(q,x)=\langle x'_1,\dots,x'_k\rangle$ and runs $T^H_U$ on the rest of the tape at $(x'_1,\dots,x'_k,x)$. On tapes $p\omega'$, every $x\in\dom D$ finds $p$ (as $p\in\dom V(\cdot,x)$, a prefix-free set), so all queries share $\omega'$, and $\nu_T(D)\ge2^{-|p|}\MH(y_{\sigma'}\Vert x_{\sigma'})$; apply (F1). For \cref{cor:hist:initial}, $\Kpred(e\text{-order})\le K(e)+c$: given $x$, enumerate $e$ until $x$ appears.
\end{proof}

To make the example of \cref{sec:hist:exploit} precise, identify $\{0,1\}^q$ with $\mathrm{GF}(2^q)$ (polynomial basis modulo the lexicographically least irreducible polynomial) and let $D_{abr}:=\{(a,r),(b,r+ab)\}$, $a\ne b$.

\begin{theorem}[A $\mathrm{GF}(2^q)$ example]\status{proved; computed}\src{history Thm 5.4, Lemma 5.5; checks c5, r3, s4}\label{thm:hist:gf}
(a) In both orders $\LH(D_{abr})\le q+c$, with $c$ independent of $q$. (b) $\LSFI(D_{abr})\ge K(a,b,r\mid q)-3q/2-1-c$; so for all but a $2^{1-j}$ fraction of triples, $\LSFI(D_{abr})\ge3q/2-j-c\ge\max_\sigma\LH(D_\sigma)+q/2-j-c'$. (c) For uniform $(a,b,r)$, $a\ne b$, $\Ex\LSFI(D_{abr})\ge3q/2-2$ for every $q\ge1$. These lower bounds are in the semimeasure convention; for $\LSFIbar$ they are open.
\end{theorem}

The key step of the proof is a character-sum bound.

\begin{lemma}[Character-sum bound]\status{proved; computed}\src{history Lemma 5.5; checks c5, s4}\label{lem:hist:charsum}
Let $F=\mathrm{GF}(2^q)$. Every partial $G:F\to F$ has $N(G):=\#\{(a,b):a\ne b,\ G(a)+G(b)=ab\}<2^q+2^{3q/2}$ (pairs where $G$ is defined).
\end{lemma}

\begin{proof}
Extending $G$ to a total function and including $a=b$ does not decrease the count. With $\chi_u(z):=(-1)^{\mathrm{Tr}(uz)}$, $\sum_u\chi_u(z)=2^q\ind[z=0]$, so in characteristic $2$
\[N\le2^{-q}\textstyle\sum_u\sum_{a,b}\chi_u(G(a))\chi_u(G(b))\chi_u(ab).\]
The term $u=0$ is $2^{2q}$. For $u\ne0$, $(\mathbf C_u)_{ab}:=\chi_u(ab)$ has $\mathbf C_u\mathbf C_u^{\mathsf T}=2^qI$, so $|v^{\mathsf T}\mathbf C_uv|\le2^{q/2}|v|^2=2^{3q/2}$ for $v_a:=\chi_u(G(a))$. Hence $N\le2^{-q}(2^{2q}+(2^q-1)2^{3q/2})<2^q+2^{3q/2}$.
\end{proof}


\begin{proof}[Proof of \cref{thm:hist:gf}]
(a) $P$ predicts $y_1$ uniformly on $\{0,1\}^{|x_1|}$ and then $y_2=y_1+x_1x_2$; in the order $(b,a)$ it predicts $(r+ab)+ba=r$. Apply \cref{cor:hist:H4}(b). (b) For fixed $\omega$, $f_\omega\supseteq D_{abr}$ iff $r=f_\omega(a)$ and $f_\omega(a)+f_\omega(b)=ab$, so $\sum_{a\ne b,r}\MF(D_{abr})=\Ex_\omega N(f_\omega|_F)<2^{3q/2+1}$ by \cref{lem:hist:charsum} (both sides count ordered pairs). So $\MF(D_{abr})2^{-3q/2-1}$ is a semimeasure on triples, lsc uniformly in $q$; apply (F2). Fewer than $2^{3q-j}$ of the at least $2^{3q-1}$ triples have $K(a,b,r\mid q)<3q-j$. (c) Jensen: $\Ex\LSFI\ge-\log\Ex\MF(D_{abr})\ge\log(2^{2q}(2^q-1))-3q/2-1\ge3q/2-2$.
\end{proof}

A direct computation (code in the accompanying repository) gives the optimum over \emph{all} laws of random functions, $\min_w\Ex[-\log w(D_{abr})]=2.708$ bits for $q=2$ and $4.348$ for $q=3$, which is a lower bound on $\Ex\LSFI$ (the one-tape version restricted to the field is a random partial function), against the Bayes history-aware loss $q$ ($\LH\le q+c$); and $\max_GN(G)=8$ for $q=2$ and $22$ for $q=3$, against the bounds $12$ and $30.6$ of \cref{lem:hist:charsum}, while for $q=4$ it lies in $[56,80]$.

\begin{theorem}[Expected gap]\status{proved}\src{history Thm 5.6}\label{thm:hist:expected}
Let $f\sim\nu$ computable and the inputs drawn by any kernel $\Delta(\cdot\mid f,h_{i-1})$, with finite alphabets up to time $n$. Then $-c\le\Ex\LSFI(D_n)-\Ex\LH\le I(f\to x_{1:n})+K(\nu)+c$, and the gap is at most $K(\nu)+c$ if the input choice does not see $f$.
\end{theorem}

The upper bound is attained up to the last input's entropy in the randomised chain, but not in general. Let $(a,b,r)$ be uniform, let $f(a)=r$, $f(b)=r+ab$ and $f=\epsilon$ elsewhere, and reveal $a$ and $b$ in a uniformly random order. Then $I(f\to x)=(q-1)+\log(2^q-1)$, while the gap $\Ex\LSFI-\Ex\LH$ is at most $q+c$. The predictor of \cref{thm:hist:gf}(a) has $\Ex L_P=q$ and its reveal law has an $O(1)$ program relative to an oracle for $f$, and $\Ex\LSFI\ge3q/2-2$: a second process on which the oracle form fails.

\begin{proof}[Proof of \cref{thm:hist:expected}]
The lower bound is \cref{lem:hist:chrono}, pathwise; the upper bound is \cref{thm:hist:directed}(b) with $P$ the predictor of $\MH$ (sub-probability conditionals), and \cref{lem:hist:memory}(a). In the $\mathrm{GF}$ process after the theorem, uniform labels give $\Ex\LSFI\le2q+c$ and $\Ex\LH$ is at least the Bayes loss $q$, so the gap is at most $q+c$. As $f$ determines $\{a,b\}$, $x_1$ has entropy $q$ and entropy $1$ given $f$, and $x_2$ is uniform on $2^q-1$ points given $h_1$ but determined by $f$ and $x_1$; so $I(f\to x)=(q-1)+\log(2^q-1)$.
\end{proof}

\begin{proposition}[An individual-sequence information bound]\status{(a) proved; (b) refuted; computed}\src{history Prop 5.7; check s1}\label{prop:hist:info}
With $I(a:b):=K(b)-K(b\mid a)$: (a) $\LSFI(D)\le\LH(D_\sigma)+I(x_\sigma:D)+c$ for every $D$ and order $\sigma$. (b) The natural strengthening with $I(x_\sigma:y_\sigma)$ in place of $I(x_\sigma:D)$ is false, even up to $O(\log n)$: on the sequences of \cref{prop:hist:strict}(i), $\LSFI(D)-\LH(D_\sigma)-I(x_\sigma:y_\sigma)\ge n-\log(n+1)-K(n)-c$ (semimeasure convention).
\end{proposition}

\begin{proof}
(a) The lookup table gives $\LSFI(D)\le K(D)+c$ (F1). For fixed $x_{1:n}$, $y_{1:n}\mapsto\MH(y_{1:n}\Vert x_{1:n})$ is a semimeasure on $Y^n$, lsc uniformly in $x_{1:n}$, so (F2) gives $K(D\mid x_\sigma)\le\LH(D_\sigma)+c$; and $K(D)=I(x_\sigma:D)+K(D\mid x_\sigma)$. (b) $\LSFI(D)\ge n-\log(n+1)-1-c$ (proof of \cref{prop:hist:strict}(i)); $\LH(D_\sigma)\le c$, as $\ind[i\text{ even}]$ knows $i$; and $I(x_\sigma:y_\sigma)\le K(y_\sigma)+c\le K(n)+c$, as $y_\sigma=0101\cdots$. With $A$ uniform, the events $\{f_\omega\supseteq D_A\}$ are disjoint, so Gibbs' inequality gives $\Ex_A\LSFI(D_A)\ge\log\binom n{\lfloor n/2\rfloor}\ge n-\log(n+1)$ exactly.
\end{proof}

